\section*{Bab \ref{ch:g5:raw-materials-and-recycling} --- Dari Bahan Baku sampai Daur Ulang}

\begin{solution}{exo:g5:raw-materials-and-recycling:1}
Alami: kayu, wol, batu. Buatan: kaca, baja, plastik.
\end{solution}

\begin{solution}{exo:g5:raw-materials-and-recycling:2}
Kertas: kayu (pohon). Kaca: pasir (bersama soda dan batu kapur). Katun:
tanaman kapas. Plastik: minyak bumi.
\end{solution}

\begin{solution}{exo:g5:raw-materials-and-recycling:3}
Batuan yang darinya sebuah logam diambil. Aluminium berasal dari
bauksit.
\end{solution}

\begin{solution}{exo:g5:raw-materials-and-recycling:4}
Stoples di tempat sampah kaca, koran bersama kertas dan karton, kaleng
bersama logam, botol bersama botol plastik.
\end{solution}

\begin{solution}{exo:g5:raw-materials-and-recycling:5}
\omterm{def:g5:raw-materials-and-recycling:raw-material}{Bahan baku} baru: bahan \omterm{def:g5:raw-materials-and-recycling:recycling}{daur ulang} dipakai menggantikan pasir, \omterm{def:g5:raw-materials-and-recycling:raw-material}{bijih},
atau minyak yang diambil dari alam.
\end{solution}

\begin{solution}{exo:g5:raw-materials-and-recycling:6}
Terbarukan: kayu (jika ditanami kembali), wol, kapas. Tidak
terbarukan: minyak bumi, \omterm{def:g5:raw-materials-and-recycling:raw-material}{bijih} besi.
\end{solution}

\begin{solution}{exo:g5:raw-materials-and-recycling:7}
Membuat kaca dari pasir adalah \omterm{def:g5:irreversible-changes:chemical-change}{perubahan kimia}: sesuatu yang baru
terbentuk, dan tidak ada jalan kembali yang sederhana.
\end{solution}

\begin{solution}{exo:g5:raw-materials-and-recycling:8}
Misalnya: mengurangi --- membeli buah tanpa bungkus plastik; menggunakan
kembali --- mengisi ulang stoples kaca; mendaur ulang --- memasukkan
kaleng minuman ke tempat sampah logam.
\end{solution}

\begin{solution}{exo:g5:raw-materials-and-recycling:9}
Setiap bahan didaur ulang dengan caranya sendiri: kepingan bahan lain
akan merusak kaca yang dilelehkan.
\end{solution}

\begin{solution}{exo:g5:raw-materials-and-recycling:10}
Seperdua puluh dari 20 satuan: kira-kira $20 \div 20 = 1$ satuan
energi.
\end{solution}

\begin{solution}{exo:g5:raw-materials-and-recycling:11}
\omterm{def:g5:raw-materials-and-recycling:recycling}{Daur ulang} tetap memerlukan energi, mesin, dan pengangkutan; benda yang
tidak pernah dibuat menghemat \omterm{def:g5:raw-materials-and-recycling:raw-material}{bahan bakunya} dan seluruh energi itu.
\end{solution}
